\documentclass[12pt]{article}
\usepackage[margin=1in]{geometry}
\usepackage{amsmath,amssymb,amsthm}
\newtheorem{theorem}{Theorem}
\newtheorem{lemma}{Lemma}
\theoremstyle{definition}
\newtheorem{assumption}{Assumption}
\theoremstyle{remark}
\newtheorem{remark}{Remark}
\newcommand{\Fnew}{F_{\mathrm{new}}}
\newcommand{\Rnew}{R_{\mathrm{new}}}
\newcommand{\ind}{\mathbf{1}}

\title{Exact finite-$K$ validity of the deployed unstudentized trajectory band\\
\large A self-contained proof of Theorem~3 (\texttt{thm:exact})}
\date{}

\begin{document}
\maketitle

This note gives an independent, self-contained proof of the exact finite-sample validity
of the deployed base band, matching its implementation exactly
(\texttt{fixed\_trajectory\_band.py}, modulation \texttt{kind="none"}: the unstudentized
``U0'' band). Nothing beyond exchangeability of the country trajectories is used---no
moment, continuity, homoskedasticity, or tail condition---which is precisely why the band
is robust to the heteroskedasticity and heavy tails that defeat its studentized variant.

\section*{Setup and notation}

Fix a horizon of $L$ recency slots and a threshold grid $\{1,\dots,T\}$. There are $K+1$
countries, indexed $c\in\{1,\dots,K,\mathrm{new}\}$, where $\mathrm{new}$ is the held-out
(target) country. Each country contributes one \emph{error trajectory}
\[
  E_c \;=\; \big(E_{c,j}(t)\big)_{j\in\{0,\dots,L-1\},\,t\in\{1,\dots,T\}} \;\in\; \mathbb{R}^{L\times T},
  \qquad
  E_{c,j}(t) \;=\; \theta_{c,j}(t) - \widehat\theta_{c,j}(t),
\]
the signed deviation of the true survey-estimate curve $\theta_{c,j}(t)$ from its predicted
center $\widehat\theta_{c,j}(t)$ at slot $j$ and threshold $t$. All countries are aligned to
exactly $L$ slots (their most recent $L$ rounds, oldest$\,\to\,$newest), so the trajectories
are index-comparable.

The deployed \emph{unstudentized} score of country $c$ is the sup-magnitude of its
trajectory,
\begin{equation}\label{eq:score}
  R_c \;=\; \|E_c\|_\infty \;=\; \max_{0\le j\le L-1}\ \max_{1\le t\le T}\ \big|E_{c,j}(t)\big|,
\end{equation}
with \emph{no} data-dependent modulation ($s\equiv 1$). Let $R_{(1)}\le\cdots\le R_{(K)}$ be
the order statistics of the $K$ calibration scores $\{R_c:c\ne\mathrm{new}\}$, put
\begin{equation}\label{eq:m}
  m \;=\; \big\lceil(1-\alpha)(K+1)\big\rceil,
\end{equation}
and define the conformal radius $\widehat q = R_{(m)}$ if $m\le K$, and $\widehat q=\infty$
otherwise. The band for every slot/threshold of the target is
\begin{equation}\label{eq:band}
  \widehat B_{j}(t) \;=\; \big[\,\widehat\theta_{\mathrm{new},j}(t)-\widehat q,\ \
  \widehat\theta_{\mathrm{new},j}(t)+\widehat q\,\big].
\end{equation}

\begin{assumption}[Trajectory exchangeability]\label{ass:exch}
The error trajectories $(E_1,\dots,E_K,E_{\mathrm{new}})$ are exchangeable random elements
of $\mathbb{R}^{L\times T}$; i.e.\ their joint law is invariant under every permutation of
the $K+1$ country indices.
\end{assumption}

\noindent
Assumption~\ref{ass:exch} is the paper's sole hypothesis. It holds whenever the country
trajectories are i.i.d.\ draws from a hyper-distribution $\Pi$ over curves and the centering
map $c\mapsto\widehat\theta_{c,\cdot}$ is applied symmetrically across countries (e.g.\ a
leave-one-country-out or full-conformal center built from a permutation-symmetric base
predictor); the round count $L$ is fixed by construction, removing trajectory-length
heterogeneity.

\section*{An exchangeable-rank coverage lemma}

\begin{lemma}\label{lem:rank}
Let $(Z_1,\dots,Z_n)$ be exchangeable real-valued random variables, and let
$Z_{(1)}\le\cdots\le Z_{(n-1)}$ be the order statistics of $(Z_1,\dots,Z_{n-1})$. For any
integer $m\in\{1,\dots,n-1\}$,
\[
  \Pr\big(Z_n \le Z_{(m)}\big)\;\ge\;\frac{m}{n},
\]
with equality when ties occur with probability zero.
\end{lemma}

\begin{proof}
Order the full sample $(Z_1,\dots,Z_n)$ by the strict total order
$(Z_i,i)\prec(Z_{i'},i')$ iff $Z_i<Z_{i'}$, or $Z_i=Z_{i'}$ and $i<i'$ (ties broken by
index). Under this order all $n$ elements are distinct, so their ranks
$\pi(1),\dots,\pi(n)$ form a permutation of $\{1,\dots,n\}$ almost surely, where $\pi(i)$ is
the position of $(Z_i,i)$. Exchangeability of $(Z_1,\dots,Z_n)$ implies the rank $\pi(n)$ of
the last coordinate is uniform on $\{1,\dots,n\}$; hence $\Pr(\pi(n)\le m)=m/n$.

If $\pi(n)\le m$, then at most $m-1$ of the first $n-1$ variables precede $Z_n$ in the
tie-broken order, so at most $m-1$ satisfy $Z_i<Z_n$ (strict), whence $Z_n\le Z_{(m)}$.
Thus $\{\pi(n)\le m\}\subseteq\{Z_n\le Z_{(m)}\}$, giving
$\Pr(Z_n\le Z_{(m)})\ge \Pr(\pi(n)\le m)=m/n$. When ties have probability zero the two
events coincide up to a null set, so equality holds.
\end{proof}

\section*{Theorem 3 and its proof}

\begin{theorem}[Exact finite-$K$ validity of the deployed base band]\label{thm:exact}
Under Assumption~\ref{ass:exch}, the band \eqref{eq:band} built from the unstudentized
scores \eqref{eq:score} with radius $\widehat q=R_{(m)}$, $m$ as in \eqref{eq:m}, satisfies
for \emph{every} $K$
\[
  \Pr\Big(\theta_{\mathrm{new},j}(t)\in\widehat B_j(t)\ \ \forall\, j,t\Big)
  \;\ge\; \frac{\big\lceil(1-\alpha)(K+1)\big\rceil}{K+1}\;\ge\;1-\alpha,
\]
distribution-free. If the calibration scores are almost surely distinct, the first
inequality is an equality (the band covers at exactly the attainable level).
\end{theorem}

\begin{proof}
If $m>K$ then $\widehat q=\infty$, the band is all of $[0,1]$ and coverage is $1\ge
\lceil(1-\alpha)(K+1)\rceil/(K+1)$; assume henceforth $m\le K$.

By \eqref{eq:band}, the target curve lies in the band at every slot and threshold iff its
deviations never exceed $\widehat q$ in magnitude:
\[
  \big\{\theta_{\mathrm{new},j}(t)\in\widehat B_j(t)\ \forall j,t\big\}
  \;=\;\Big\{\max_{j,t}\big|E_{\mathrm{new},j}(t)\big|\le\widehat q\Big\}
  \;=\;\big\{\Rnew\le R_{(m)}\big\},
\]
using the definition \eqref{eq:score} of the score. Now the scores are obtained by applying
the \emph{same} fixed measurable map $\varphi(E)=\|E\|_\infty$ to each trajectory,
$R_c=\varphi(E_c)$. Because $(E_1,\dots,E_K,E_{\mathrm{new}})$ are exchangeable
(Assumption~\ref{ass:exch}) and $\varphi$ does not depend on $c$, the scores
$(R_1,\dots,R_K,\Rnew)$ are exchangeable. Applying Lemma~\ref{lem:rank} with $n=K+1$,
$Z_n=\Rnew$, and $m$ as in \eqref{eq:m},
\[
  \Pr\big(\Rnew\le R_{(m)}\big)\;\ge\;\frac{m}{K+1}
  \;=\;\frac{\lceil(1-\alpha)(K+1)\rceil}{K+1}\;\ge\;\frac{(1-\alpha)(K+1)}{K+1}\;=\;1-\alpha,
\]
and the equality claim under a.s.-distinct scores is the equality clause of
Lemma~\ref{lem:rank}. \qedhere
\end{proof}

\begin{remark}[Why studentization forfeits exactness]
Replacing \eqref{eq:score} by a studentized score $R_c=\max_{j,t}|E_{c,j}(t)|/\widehat s_j(t)$
with $\widehat s_j(t)$ an in-sample spread estimated from the calibration errors destroys the
symmetry used above: each \emph{calibration} score is divided by an $\widehat s$ computed from
a set that \emph{includes} its own curve, whereas the \emph{target} score is divided by an
$\widehat s$ computed from a set that \emph{excludes} the target. The map producing $R_c$ is
then no longer a common function of $E_c$ alone, the scores are no longer exchangeable, and a
self-inclusion asymmetry of order $O(1/K)$ enters: the studentized band attains $1-\alpha$
only as $K\to\infty$. The deployed band therefore fixes $\widehat s\equiv 1$; on the
near-homoskedastic real curves the resulting width cost is $\le 5\%$.
\end{remark}

\begin{remark}[Scale-freeness and the contract test]
Since no scale is estimated, multiplying every error by a constant $\lambda>0$ multiplies
both $\Rnew$ and $\widehat q$ by $\lambda$ and leaves the covered event
$\{\Rnew\le\widehat q\}$ unchanged; heteroskedasticity and heavy tails thus cannot break the
rank symmetry. The guarantee is verified numerically in
\texttt{test\_unstudentized\_coverage\_at\_attainable\_level} (K=30, heteroskedastic
$t_3$-tailed errors): the empirical coverage sits at the attainable level
$\lceil(1-\alpha)(K+1)\rceil/(K+1)$ within Monte Carlo error, and
\texttt{test\_unstudentized\_score\_is\_scale\_free} checks the scale invariance above.
\end{remark}

\end{document}
